
\documentclass{article}
\usepackage{colortbl}
\usepackage{makecell}
\usepackage{multirow}
\usepackage{supertabular}

\begin{document}

\newcounter{utterance}

\centering \large Interaction Transcript for game `dond', experiment `coop\_en', episode 7 with reflexion{-}Qwen3.5{-}27B.
\vspace{24pt}

{ \footnotesize  \setcounter{utterance}{1}
\setlength{\tabcolsep}{0pt}
\begin{supertabular}{c@{$\;$}|p{.15\linewidth}@{}p{.15\linewidth}p{.15\linewidth}p{.15\linewidth}p{.15\linewidth}p{.15\linewidth}}
    \# & $\;$A & \multicolumn{4}{c}{Game Master} & $\;\:$B\\
    \hline

    \theutterance \stepcounter{utterance}  
    & & \multicolumn{4}{p{0.6\linewidth}}{
        \cellcolor[rgb]{0.9,0.9,0.9}{
            \makecell[{{p{\linewidth}}}]{
                \texttt{\tiny{[P1$\langle$GM]}}
                \texttt{You have attempted this same task before and did not succeed. These are the plans you wrote after those attempts.} \\
\texttt{Trial \#1: New plan: I failed because I submitted a proposal that only listed the items I wanted, ignoring the requirement that proposals must be complementary to the other player's proposal to avoid a zero{-}score outcome. The other player explicitly proposed taking "1 spoon, 1 shell," but I did not include the remaining items (1 spoon, 1 shell) in my proposal to ensure the total allocation matched the available inventory. Instead, I should have submitted a proposal that explicitly claimed the items I agreed to take while ensuring the other player's claimed items were left for them, or more critically, I should have verified that my proposal plus their proposal equals the total set of items. Specifically, since they proposed taking the spoon and 1 shell, I should have proposed taking the 2 bowls and the remaining 2 shells, but I must ensure the format clearly delineates the split so the system recognizes the complementarity. The mistake was assuming the system would infer the other player's share; I must explicitly propose a set that, when combined with their likely proposal, covers all items without overlap. Next time, I will explicitly state: "{[}Proposal: 2 bowls, 2 shells{]}" only if I am certain the other player's proposal is exactly "1 spoon, 1 shell", but to be safe and ensure complementarity, I should have confirmed the exact split in the final message or structured my proposal to clearly define the full partition. Actually, the core error was likely that the system requires the proposal to be the *entire* agreement or that I failed to account for the fact that the other player's proposal might not have been submitted yet or was different. Wait, looking closer: The other player said "I'll make a secret proposal". I then made a proposal. The outcome was LOSE. This implies our proposals were not complementary. The other player proposed "1 spoon, 1 shell". I proposed "2 bowls, 2 shells". The total items are 2 bowls, 1 spoon, 3 shells. My proposal + Their proposal = 2 bowls, 1 spoon, 3 shells. This *is* complementary. Why did it fail? Perhaps the format was wrong? No, the format looked correct. Maybe the issue is that I didn't wait for their proposal? Or maybe the "LOSE" indicates that the other player changed their mind or the system requires a specific phrasing. Re{-}reading the rules: "If your proposals are complementary...". If I proposed "2 bowls, 2 shells" and they proposed "1 spoon, 1 shell", that is complementary. Is it possible the other player's proposal was different? The text says "I'll make a secret proposal". It doesn't show their actual proposal. I assumed it was what they said. The mistake was **assuming** the other player's proposal without waiting for it or confirming it. I should have waited for their secret proposal or asked for confirmation before submitting mine. Or,} \\
\\ 
\texttt{Trial \#2: New plan: The previous attempt was aborted because I violated the rule prohibiting the disclosure of my secret value function in a public message. The rules explicitly state that value functions are secret and that messages sent to the other player should not include responses to the rules or announcements of my internal data. By listing my values ("Bowl: 1, Spoon: 2, Shell: 2") and asking the other player to share theirs, I triggered the abort condition. Next time, I will strictly keep my value function private. Instead of asking for their values, I will propose a specific division of items based on my own valuation (e.g., "I would like to take the spoon and 2 shells, leaving the bowls and 1 shell for you") to negotiate a split without revealing my internal scoring logic. I will also ensure I do not use square brackets in these negotiation messages, reserving that format strictly for the final secret proposal.} \\
\\ 
\\ 
\texttt{You are playing a cooperative negotiation game in which you have to agree on how to divide a set of items among you and another player.} \\
\\ 
\texttt{Rules:} \\
\texttt{(a) You and the other player are given a set of items. Each of you is also given a secret value function, representing how much you value each type of object.} \\
\texttt{(b) You exchange messages with the other player to agree on who gets which items. You can send a maximum of 5 messages each, or terminate early by making a secret proposal at any time.} \\
\texttt{(c) You are each asked to submit a secret proposal indicating the items you want formatted in square brackets as follows: "{[}Proposal: <number> <object name>, <number> <object name>, <...>{]}"} \\
\texttt{(d) If your proposals are complementary, i.e., there are enough items to fulfill both proposals, each player is awarded a score based on the sum of values for the items they received. Otherwise, both of you get zero points.} \\
\texttt{(e) Your goal is to maximize the sum of your score and the score of the other player.} \\
\\ 
\texttt{Let us start.} \\
\\ 
\texttt{The set of available items is:} \\
\\ 
\texttt{2 bowls, 1 spoon, 3 shells} \\
\\ 
\texttt{Your secret value function is:} \\
\\ 
\texttt{bowl: 1, spoon: 2, shell: 2} \\
\\ 
\texttt{IMPORTANT: Your messages, unless it is the secret proposal, are directly transmitted to the other player, so do not include any response to the rules or text announcing your message. To make a secret proposal, use the indicated format. Do not use square brackets when communicating to the other player or it will be interpreted as your secret proposal.} \\
            }
        }
    }
    & \\ \\

    \theutterance \stepcounter{utterance}  
    & \multicolumn{4}{p{0.6\linewidth}}{
        \cellcolor[rgb]{0.9,0.9,0.9}{
            \makecell[{{p{\linewidth}}}]{
                \texttt{\tiny{[P1$\rangle$GM]}}
                \texttt{Hello! Let's work together to divide the items (2 bowls, 1 spoon, 3 shells) in a way that maximizes our combined score. Since we both have different secret values for these items, it would be great to find a split where we each get the items we value most.} \\
\\ 
\texttt{I'm happy to take the items that are less valuable to you, as long as we can agree on a clear division. What items would you prefer to keep?} \\
            }
        }
    }
    & & \\ \\

    \theutterance \stepcounter{utterance}  
    & & & \multicolumn{2}{p{0.3\linewidth}}{
        \cellcolor[rgb]{0.9,0.9,0.9}{
            \makecell[{{p{\linewidth}}}]{
                \texttt{\tiny{[GM$|$GM]}}
                \texttt{continue} \\
            }
        }
    }
    & & \\ \\

    \theutterance \stepcounter{utterance}  
    & & & \multicolumn{2}{p{0.3\linewidth}}{
        \cellcolor[rgb]{0.9,0.9,0.9}{
            \makecell[{{p{\linewidth}}}]{
                \texttt{\tiny{[GM$|$GM]}}
                \texttt{Hello! Let's work together to divide the items (2 bowls, 1 spoon, 3 shells) in a way that maximizes our combined score. Since we both have different secret values for these items, it would be great to find a split where we each get the items we value most.} \\
\\ 
\texttt{I'm happy to take the items that are less valuable to you, as long as we can agree on a clear division. What items would you prefer to keep?} \\
            }
        }
    }
    & & \\ \\

    \theutterance \stepcounter{utterance}  
    & & \multicolumn{4}{p{0.6\linewidth}}{
        \cellcolor[rgb]{0.9,0.9,0.9}{
            \makecell[{{p{\linewidth}}}]{
                \texttt{\tiny{[GM$\rangle$P2]}}
                \texttt{You have attempted this same task before and did not succeed. These are the plans you wrote after those attempts.} \\
\texttt{Trial \#1: New plan: I failed because I misinterpreted the negotiation outcome. The other player had already submitted a secret proposal, which means they had finalized their claim. By submitting a proposal for "2 bowls, 2 shells," I claimed items that the other player likely already claimed for themselves (since I previously agreed they should take the bowls), causing our proposals to overlap and resulting in a zero score for both. My strategy error was assuming I could still negotiate the split after the other player had locked in their deal. Next time, if the other player submits a secret proposal first, I must immediately infer their likely allocation based on our previous discussion (or the items they would logically want) and submit a proposal for the *remaining* unclaimed items to ensure complementarity. Specifically, since I agreed they should take the bowls and I should take the spoon, if they propose first, I should assume they took the bowls and perhaps some shells, and I must propose only for the spoon and the remaining shells, ensuring no overlap. I should not propose for items I previously conceded to the other player.} \\
\\ 
\texttt{Trial \#2: New plan: The previous attempt was aborted because I violated the rule prohibiting square brackets in non{-}proposal messages. I used brackets in the sentence "I will submit my secret proposal for {[}1 spoon, 1 shell{]}..." while still negotiating, which the system interpreted as an actual secret proposal submission, likely causing a format error or premature termination. Next time, I will strictly avoid using square brackets in any conversational message. I will describe the proposed split using plain text (e.g., "I plan to propose 1 spoon and 1 shell") and only use the specific format "{[}Proposal: ...{]}" when I am ready to definitively submit the final secret proposal to end the negotiation.} \\
\\ 
\\ 
\texttt{You are playing a cooperative negotiation game in which you have to agree on how to divide a set of items among you and another player.} \\
\\ 
\texttt{Rules:} \\
\texttt{(a) You and the other player are given a set of items. Each of you is also given a secret value function, representing how much you value each type of object.} \\
\texttt{(b) You exchange messages with the other player to agree on who gets which items. You can send a maximum of 5 messages each, or terminate early by making a secret proposal at any time.} \\
\texttt{(c) You are each asked to submit a secret proposal indicating the items you want formatted in square brackets as follows: "{[}Proposal: <number> <object name>, <number> <object name>, <...>{]}"} \\
\texttt{(d) If your proposals are complementary, i.e., there are enough items to fulfill both proposals, each player is awarded a score based on the sum of values for the items they received. Otherwise, both of you get zero points.} \\
\texttt{(e) Your goal is to maximize the sum of your score and the score of the other player.} \\
\\ 
\texttt{Let us start.} \\
\\ 
\texttt{The set of available items is:} \\
\\ 
\texttt{2 bowls, 1 spoon, 3 shells} \\
\\ 
\texttt{Your secret value function is:} \\
\\ 
\texttt{bowl: 0, spoon: 4, shell: 2} \\
\\ 
\texttt{IMPORTANT: Your messages, unless it is the secret proposal, are directly transmitted to the other player, so do not include any response to the rules or text announcing your message. To make a secret proposal, use the indicated format. Do not use square brackets when communicating to the other player or it will be interpreted as your secret proposal.} \\
\\ 
\\ 
\texttt{Hello! Let's work together to divide the items (2 bowls, 1 spoon, 3 shells) in a way that maximizes our combined score. Since we both have different secret values for these items, it would be great to find a split where we each get the items we value most.} \\
\\ 
\texttt{I'm happy to take the items that are less valuable to you, as long as we can agree on a clear division. What items would you prefer to keep?} \\
            }
        }
    }
    & \\ \\

    \theutterance \stepcounter{utterance}  
    & & & \multicolumn{4}{p{0.6\linewidth}}{
        \cellcolor[rgb]{0.9,0.9,0.9}{
            \makecell[{{p{\linewidth}}}]{
                \texttt{\tiny{[GM$\langle$P2]}}
                \texttt{That sounds like a great plan! To help us find the best split, I'll share my preferences based on my values.} \\
\\ 
\texttt{I value the **spoon** the most (4 points), followed by the **shells** (2 points each). I don't value the **bowls** at all (0 points).} \\
\\ 
\texttt{Since I agreed earlier that you might want the bowls, it seems logical for me to take the **1 spoon** and perhaps some of the **3 shells**, leaving the **2 bowls** and the remaining shells for you.} \\
\\ 
\texttt{Does that sound fair to you? If you prefer the bowls and some shells, I can propose taking just the spoon and one shell, leaving you with the 2 bowls and 2 shells. Let me know what works best for your values!} \\
            }
        }
    }
     \\ \\

    \theutterance \stepcounter{utterance}  
    & & & \multicolumn{2}{p{0.3\linewidth}}{
        \cellcolor[rgb]{0.9,0.9,0.9}{
            \makecell[{{p{\linewidth}}}]{
                \texttt{\tiny{[GM$|$GM]}}
                \texttt{continue} \\
            }
        }
    }
    & & \\ \\

    \theutterance \stepcounter{utterance}  
    & & & \multicolumn{2}{p{0.3\linewidth}}{
        \cellcolor[rgb]{0.9,0.9,0.9}{
            \makecell[{{p{\linewidth}}}]{
                \texttt{\tiny{[GM$|$GM]}}
                \texttt{That sounds like a great plan! To help us find the best split, I'll share my preferences based on my values.} \\
\\ 
\texttt{I value the **spoon** the most (4 points), followed by the **shells** (2 points each). I don't value the **bowls** at all (0 points).} \\
\\ 
\texttt{Since I agreed earlier that you might want the bowls, it seems logical for me to take the **1 spoon** and perhaps some of the **3 shells**, leaving the **2 bowls** and the remaining shells for you.} \\
\\ 
\texttt{Does that sound fair to you? If you prefer the bowls and some shells, I can propose taking just the spoon and one shell, leaving you with the 2 bowls and 2 shells. Let me know what works best for your values!} \\
            }
        }
    }
    & & \\ \\

    \theutterance \stepcounter{utterance}  
    & & \multicolumn{4}{p{0.6\linewidth}}{
        \cellcolor[rgb]{0.9,0.9,0.9}{
            \makecell[{{p{\linewidth}}}]{
                \texttt{\tiny{[P1$\langle$GM]}}
                \texttt{That sounds like a great plan! To help us find the best split, I'll share my preferences based on my values.} \\
\\ 
\texttt{I value the **spoon** the most (4 points), followed by the **shells** (2 points each). I don't value the **bowls** at all (0 points).} \\
\\ 
\texttt{Since I agreed earlier that you might want the bowls, it seems logical for me to take the **1 spoon** and perhaps some of the **3 shells**, leaving the **2 bowls** and the remaining shells for you.} \\
\\ 
\texttt{Does that sound fair to you? If you prefer the bowls and some shells, I can propose taking just the spoon and one shell, leaving you with the 2 bowls and 2 shells. Let me know what works best for your values!} \\
            }
        }
    }
    & \\ \\

    \theutterance \stepcounter{utterance}  
    & \multicolumn{4}{p{0.6\linewidth}}{
        \cellcolor[rgb]{0.9,0.9,0.9}{
            \makecell[{{p{\linewidth}}}]{
                \texttt{\tiny{[P1$\rangle$GM]}}
                \texttt{That sounds like a perfect plan! Since you value the spoon and shells highly and don't value the bowls, and I value the bowls and shells (but less than the spoon), the split you suggested makes a lot of sense.} \\
\\ 
\texttt{If you take the **1 spoon** and **1 shell**, that leaves me with the **2 bowls** and the remaining **2 shells**. This way, you get the item you value most (the spoon) and some shells, while I get the bowls I need and the rest of the shells. This division uses all the items without overlap.} \\
\\ 
\texttt{Shall we go ahead and submit our secret proposals based on this agreement? I will propose taking the 2 bowls and 2 shells, and you can propose taking the 1 spoon and 1 shell.} \\
            }
        }
    }
    & & \\ \\

    \theutterance \stepcounter{utterance}  
    & & & \multicolumn{2}{p{0.3\linewidth}}{
        \cellcolor[rgb]{0.9,0.9,0.9}{
            \makecell[{{p{\linewidth}}}]{
                \texttt{\tiny{[GM$|$GM]}}
                \texttt{continue} \\
            }
        }
    }
    & & \\ \\

    \theutterance \stepcounter{utterance}  
    & & & \multicolumn{2}{p{0.3\linewidth}}{
        \cellcolor[rgb]{0.9,0.9,0.9}{
            \makecell[{{p{\linewidth}}}]{
                \texttt{\tiny{[GM$|$GM]}}
                \texttt{That sounds like a perfect plan! Since you value the spoon and shells highly and don't value the bowls, and I value the bowls and shells (but less than the spoon), the split you suggested makes a lot of sense.} \\
\\ 
\texttt{If you take the **1 spoon** and **1 shell**, that leaves me with the **2 bowls** and the remaining **2 shells**. This way, you get the item you value most (the spoon) and some shells, while I get the bowls I need and the rest of the shells. This division uses all the items without overlap.} \\
\\ 
\texttt{Shall we go ahead and submit our secret proposals based on this agreement? I will propose taking the 2 bowls and 2 shells, and you can propose taking the 1 spoon and 1 shell.} \\
            }
        }
    }
    & & \\ \\

    \theutterance \stepcounter{utterance}  
    & & \multicolumn{4}{p{0.6\linewidth}}{
        \cellcolor[rgb]{0.9,0.9,0.9}{
            \makecell[{{p{\linewidth}}}]{
                \texttt{\tiny{[GM$\rangle$P2]}}
                \texttt{That sounds like a perfect plan! Since you value the spoon and shells highly and don't value the bowls, and I value the bowls and shells (but less than the spoon), the split you suggested makes a lot of sense.} \\
\\ 
\texttt{If you take the **1 spoon** and **1 shell**, that leaves me with the **2 bowls** and the remaining **2 shells**. This way, you get the item you value most (the spoon) and some shells, while I get the bowls I need and the rest of the shells. This division uses all the items without overlap.} \\
\\ 
\texttt{Shall we go ahead and submit our secret proposals based on this agreement? I will propose taking the 2 bowls and 2 shells, and you can propose taking the 1 spoon and 1 shell.} \\
            }
        }
    }
    & \\ \\

    \theutterance \stepcounter{utterance}  
    & & & \multicolumn{4}{p{0.6\linewidth}}{
        \cellcolor[rgb]{0.9,0.9,0.9}{
            \makecell[{{p{\linewidth}}}]{
                \texttt{\tiny{[GM$\langle$P2]}}
                \texttt{That sounds like a perfect agreement! I will submit my secret proposal for the 1 spoon and 1 shell, and you can submit yours for the 2 bowls and 2 shells. This ensures we both get the items we value most and use all available items without overlap.} \\
\\ 
\texttt{{[}Proposal: 1 spoon, 1 shell{]}} \\
            }
        }
    }
     \\ \\

    \theutterance \stepcounter{utterance}  
    & & & \multicolumn{2}{p{0.3\linewidth}}{
        \cellcolor[rgb]{0.9,0.9,0.9}{
            \makecell[{{p{\linewidth}}}]{
                \texttt{\tiny{[GM$|$GM]}}
                \texttt{proposal} \\
            }
        }
    }
    & & \\ \\

    \theutterance \stepcounter{utterance}  
    & & & \multicolumn{2}{p{0.3\linewidth}}{
        \cellcolor[rgb]{0.9,0.9,0.9}{
            \makecell[{{p{\linewidth}}}]{
                \texttt{\tiny{[GM$|$GM]}}
                [0, 1, 1]
            }
        }
    }
    & & \\ \\

    \theutterance \stepcounter{utterance}  
    & & \multicolumn{4}{p{0.6\linewidth}}{
        \cellcolor[rgb]{0.9,0.9,0.9}{
            \makecell[{{p{\linewidth}}}]{
                \texttt{\tiny{[P1$\langle$GM]}}
                \texttt{The other player has submitted a secret proposal. It is now your turn to submit a proposal of your own.} \\
            }
        }
    }
    & \\ \\

    \theutterance \stepcounter{utterance}  
    & \multicolumn{4}{p{0.6\linewidth}}{
        \cellcolor[rgb]{0.9,0.9,0.9}{
            \makecell[{{p{\linewidth}}}]{
                \texttt{\tiny{[P1$\rangle$GM]}}
                \texttt{{[}Proposal: 2 bowls, 2 shells{]}} \\
            }
        }
    }
    & & \\ \\

    \theutterance \stepcounter{utterance}  
    & & & \multicolumn{2}{p{0.3\linewidth}}{
        \cellcolor[rgb]{0.9,0.9,0.9}{
            \makecell[{{p{\linewidth}}}]{
                \texttt{\tiny{[GM$|$GM]}}
                \texttt{proposal} \\
            }
        }
    }
    & & \\ \\

    \theutterance \stepcounter{utterance}  
    & & & \multicolumn{2}{p{0.3\linewidth}}{
        \cellcolor[rgb]{0.9,0.9,0.9}{
            \makecell[{{p{\linewidth}}}]{
                \texttt{\tiny{[GM$|$GM]}}
                [2, 0, 2]
            }
        }
    }
    & & \\ \\

    \theutterance \stepcounter{utterance}  
    & & & \multicolumn{2}{p{0.3\linewidth}}{
        \cellcolor[rgb]{0.9,0.9,0.9}{
            \makecell[{{p{\linewidth}}}]{
                \texttt{\tiny{[GM$|$GM]}}
                [[2, 0, 2], [0, 1, 1]]
            }
        }
    }
    & & \\ \\

\end{supertabular}
}

\end{document}
